\chapter[Discusión]{Discusión general}
\label{cp:discusion}
% Extensión estimada: 8--12 páginas.
% Fuentes: sección "Discusión y conclusiones" del documento de experimentación (incluye "Amenazas a la
% validez" y "Limitaciones") y la comparación de hardware del antiguo capítulo de costo computacional.
%
% Feedback del tutor (14/09/2026): cada capítulo de resultados cierra con su propia discusión; este capítulo
% queda para lo transversal, lo que ninguno de los dos podría cerrar por sí solo. Por eso se fusionan aquí
% "Comparación con trabajos relacionados" dentro de las secciones temáticas y "Amenazas a la validez" con
% "Limitaciones".

\pendiente{Este capítulo interpreta en conjunto los resultados de los dos capítulos anteriores y los ubica
frente a la literatura, sin repetir sus tablas.}

\section{Respuesta a las preguntas de investigación}
\label{sec:disc_preguntas}
\pendiente{Se responden la PI1 y la PI2 con la evidencia de los capítulos de resultados y se resume el estado
final de cada hipótesis.}

\section{Especialización de dominio frente a modalidad}
\label{sec:disc_especializacion}
\pendiente{Se discute por qué el factor determinante no es que el modelo sea visión--lenguaje, sino que esté
especializado en dermatología, y cómo se ubica esa conclusión respecto de los trabajos revisados en la
Sección~\ref{sec:estado_arte}.}

\section{División de trabajo entre clasificador y generador}
\label{sec:disc_division}
\pendiente{Se discute que el clasificador aporta el diagnóstico y el VLM la descripción, y qué implica esto
para el diseño de sistemas de generación de reportes médicos.}

\section{Fidelidad diagnóstica y validez clínica}
\label{sec:disc_validez_clinica}
\pendiente{Se delimita qué establecen y qué no las métricas automáticas, y el papel de la incertidumbre y la
derivación en un sistema de apoyo.}

\section{Lecciones metodológicas}
\label{sec:disc_lecciones}
\pendiente{Se discuten la fuga de información por lesión en HAM10000 y su efecto sobre los resultados
publicados en el área, la importancia de seleccionar exclusivamente en validación, la inestabilidad entre
semillas de las métricas de las clases minoritarias y el sesgo de las métricas con referencia al comparar
estilos de reporte.}

\section{Viabilidad en hardware accesible}
\label{sec:disc_hardware}
\pendiente{Se relacionan el desempeño alcanzado y el costo medido en la Sección~\ref{sec:res_costo} con los
requerimientos de cómputo de los trabajos del estado del arte, para discutir qué parte de la brecha con el
SOTA se explica por escala de modelo y de datos y qué parte no.}

\section{Limitaciones y amenazas a la validez}
\label{sec:disc_limitaciones}
\pendiente{Se analizan las amenazas a la validez interna, externa, de constructo y de conclusión estadística,
junto con las limitaciones del trabajo: escala de los modelos, costo, ausencia de reportes redactados por
especialistas, generalización fuera de distribución, anclajes y generadores no evaluados, y generación de
reportes sobre el test original.

Para cada hipótesis exploratoria (H3 a H6) se indica si hubo evidencia independiente de los datos que la
motivaron (HIBA, la partición disjunta por lesión, semillas nuevas) o si se contrastó sobre el mismo
material, en cuyo caso su respaldo queda como exploratorio. La Sección~\ref{sec:intro_preguntas} remite
aquí.}
